% GUID v6, Primer P9 (label box:particle): particle physics. Background only; it states no result of the paper.
% Place it before "The spin ladder" (marker %PRIMER:particle in part1.tex): % GUID v6, Primer P9 (label box:particle): particle physics. Background only; it states no result of the paper.
% Place it before "The spin ladder" (marker %PRIMER:particle in part1.tex): % GUID v6, Primer P9 (label box:particle): particle physics. Background only; it states no result of the paper.
% Place it before "The spin ladder" (marker %PRIMER:particle in part1.tex): % GUID v6, Primer P9 (label box:particle): particle physics. Background only; it states no result of the paper.
% Place it before "The spin ladder" (marker %PRIMER:particle in part1.tex): \input{primers/p9-particle}
\begin{primer}{Particle physics}{box:particle}
\setlength{\parskip}{3pt plus 1pt}%
\emph{The Higgs mechanism.} The Standard Model is a gauge theory of $SU(3)\times SU(2)\times U(1)$: colour for the strong force, weak isospin
and hypercharge for the electroweak force~\cite{glashow,weinberg67,salam}. The Higgs field is a doublet whose potential has its minimum
away from zero, at a vacuum value $v=246$\,GeV~\cite{higgs64,be64}. In that vacuum the W and Z bosons acquire masses proportional to $v$,
while one combination of the electroweak fields, the photon, stays massless. Quarks and charged leptons get their masses from
\emph{Yukawa couplings} $y_f$ to the doublet, $m_f=y_fv/\sqrt2$. The Higgs boson therefore couples to each fermion with strength $m_f/v$
and to the W and Z with $2m_V^2/v$.

\emph{The $\kappa$ framework.} Measurements of Higgs production and decay are summarised by \emph{coupling modifiers} $\kappa_i$, which
scale each Standard Model coupling, so that a rate involving particle $i$ scales as $\kappa_i^2$; the Standard Model has every
$\kappa_i=1$~\cite{david2012}.

\emph{Schemes and running.} A coupling or mass in quantum field theory is defined only together with a \emph{renormalisation scheme} and
an energy scale $\mu$. In the $\overline{\text{MS}}$ scheme only the divergent part of each loop is subtracted, and parameters run with
$\mu$. The \emph{on-shell} scheme instead fixes masses at the poles of the propagators and $\alpha$ at zero momentum transfer. The same
quantity can differ between the two by more than its measurement error, so every comparison states its scheme~\cite{erler2017,navas2024}.
The dependence on $\mu$ obeys the renormalisation-group equations (RGEs). At one loop
$\dd\alpha_i^{-1}/\dd\ln\mu=-b_i/2\pi$, with $(b_1,b_2,b_3)=(41/10,-19/6,-7)$ in the Standard Model (hypercharge normalised as in grand
unification). The strong coupling therefore weakens at high energy, and the hypercharge coupling grows. Precision work runs at two or
three loops~\cite{buttazzo2013}.

\emph{Dirac or Majorana.} A \emph{Dirac} mass joins a left-handed neutrino to a right-handed partner and conserves lepton number $L$. A
\emph{Majorana} mass joins the neutrino to its own antiparticle and changes $L$ by two. The lowest-dimension operator of Standard Model
fields that gives neutrinos mass is Weinberg's dimension-five operator, $(c/\Lambda)\,(LH)(LH)$, with $L$ the lepton doublet; in the Higgs
vacuum it gives a Majorana mass of order $cv^2/\Lambda$, small when the new scale $\Lambda$ is large~\cite{weinberg1979}. Majorana masses
can therefore also come from the Higgs vacuum.

\emph{Neutrinoless double-beta decay} ($0\nu\beta\beta$) turns two neutrons of a nucleus into two protons and two electrons with no
neutrinos. It changes $L$ by two and can occur only if neutrinos have a Majorana mass. Through light-neutrino exchange its rate is
proportional to $m_{\beta\beta}^2$, with the effective mass $m_{\beta\beta}=|\sum_iU_{ei}^2m_i|$ set by the mixing matrix $U$ and the
masses $m_i$~\cite{navas2024}. Xenon and germanium experiments search for it~\cite{kamlandzen2025,legend2025}.

\emph{Grand unification.} A grand unified theory (GUT) embeds the three gauge groups in one simple group, the smallest being
$SU(5)$~\cite{georgiglashow}, with a single coupling above a scale $M_X$. With the Standard Model's particles alone, the three running
couplings do not meet at one point~\cite{amaldi}. A GUT has two generic consequences. Its heavy gauge bosons let protons decay, with a
lifetime of order $M_X^4/(\alpha_G^2m_p^5)$; Super-Kamiokande's bound on $p\to e^+\pi^0$ excludes the minimal $SU(5)$
model~\cite{superk}. And its breaking leaves magnetic monopoles of mass near $M_X/\alpha_G$, about one per causal volume, which would
overfill the universe unless inflation dilutes them~\cite{preskill,guth}.
\end{primer}

\begin{primer}{Particle physics}{box:particle}
\setlength{\parskip}{3pt plus 1pt}%
\emph{The Higgs mechanism.} The Standard Model is a gauge theory of $SU(3)\times SU(2)\times U(1)$: colour for the strong force, weak isospin
and hypercharge for the electroweak force~\cite{glashow,weinberg67,salam}. The Higgs field is a doublet whose potential has its minimum
away from zero, at a vacuum value $v=246$\,GeV~\cite{higgs64,be64}. In that vacuum the W and Z bosons acquire masses proportional to $v$,
while one combination of the electroweak fields, the photon, stays massless. Quarks and charged leptons get their masses from
\emph{Yukawa couplings} $y_f$ to the doublet, $m_f=y_fv/\sqrt2$. The Higgs boson therefore couples to each fermion with strength $m_f/v$
and to the W and Z with $2m_V^2/v$.

\emph{The $\kappa$ framework.} Measurements of Higgs production and decay are summarised by \emph{coupling modifiers} $\kappa_i$, which
scale each Standard Model coupling, so that a rate involving particle $i$ scales as $\kappa_i^2$; the Standard Model has every
$\kappa_i=1$~\cite{david2012}.

\emph{Schemes and running.} A coupling or mass in quantum field theory is defined only together with a \emph{renormalisation scheme} and
an energy scale $\mu$. In the $\overline{\text{MS}}$ scheme only the divergent part of each loop is subtracted, and parameters run with
$\mu$. The \emph{on-shell} scheme instead fixes masses at the poles of the propagators and $\alpha$ at zero momentum transfer. The same
quantity can differ between the two by more than its measurement error, so every comparison states its scheme~\cite{erler2017,navas2024}.
The dependence on $\mu$ obeys the renormalisation-group equations (RGEs). At one loop
$\dd\alpha_i^{-1}/\dd\ln\mu=-b_i/2\pi$, with $(b_1,b_2,b_3)=(41/10,-19/6,-7)$ in the Standard Model (hypercharge normalised as in grand
unification). The strong coupling therefore weakens at high energy, and the hypercharge coupling grows. Precision work runs at two or
three loops~\cite{buttazzo2013}.

\emph{Dirac or Majorana.} A \emph{Dirac} mass joins a left-handed neutrino to a right-handed partner and conserves lepton number $L$. A
\emph{Majorana} mass joins the neutrino to its own antiparticle and changes $L$ by two. The lowest-dimension operator of Standard Model
fields that gives neutrinos mass is Weinberg's dimension-five operator, $(c/\Lambda)\,(LH)(LH)$, with $L$ the lepton doublet; in the Higgs
vacuum it gives a Majorana mass of order $cv^2/\Lambda$, small when the new scale $\Lambda$ is large~\cite{weinberg1979}. Majorana masses
can therefore also come from the Higgs vacuum.

\emph{Neutrinoless double-beta decay} ($0\nu\beta\beta$) turns two neutrons of a nucleus into two protons and two electrons with no
neutrinos. It changes $L$ by two and can occur only if neutrinos have a Majorana mass. Through light-neutrino exchange its rate is
proportional to $m_{\beta\beta}^2$, with the effective mass $m_{\beta\beta}=|\sum_iU_{ei}^2m_i|$ set by the mixing matrix $U$ and the
masses $m_i$~\cite{navas2024}. Xenon and germanium experiments search for it~\cite{kamlandzen2025,legend2025}.

\emph{Grand unification.} A grand unified theory (GUT) embeds the three gauge groups in one simple group, the smallest being
$SU(5)$~\cite{georgiglashow}, with a single coupling above a scale $M_X$. With the Standard Model's particles alone, the three running
couplings do not meet at one point~\cite{amaldi}. A GUT has two generic consequences. Its heavy gauge bosons let protons decay, with a
lifetime of order $M_X^4/(\alpha_G^2m_p^5)$; Super-Kamiokande's bound on $p\to e^+\pi^0$ excludes the minimal $SU(5)$
model~\cite{superk}. And its breaking leaves magnetic monopoles of mass near $M_X/\alpha_G$, about one per causal volume, which would
overfill the universe unless inflation dilutes them~\cite{preskill,guth}.
\end{primer}

\begin{primer}{Particle physics}{box:particle}
\setlength{\parskip}{3pt plus 1pt}%
\emph{The Higgs mechanism.} The Standard Model is a gauge theory of $SU(3)\times SU(2)\times U(1)$: colour for the strong force, weak isospin
and hypercharge for the electroweak force~\cite{glashow,weinberg67,salam}. The Higgs field is a doublet whose potential has its minimum
away from zero, at a vacuum value $v=246$\,GeV~\cite{higgs64,be64}. In that vacuum the W and Z bosons acquire masses proportional to $v$,
while one combination of the electroweak fields, the photon, stays massless. Quarks and charged leptons get their masses from
\emph{Yukawa couplings} $y_f$ to the doublet, $m_f=y_fv/\sqrt2$. The Higgs boson therefore couples to each fermion with strength $m_f/v$
and to the W and Z with $2m_V^2/v$.

\emph{The $\kappa$ framework.} Measurements of Higgs production and decay are summarised by \emph{coupling modifiers} $\kappa_i$, which
scale each Standard Model coupling, so that a rate involving particle $i$ scales as $\kappa_i^2$; the Standard Model has every
$\kappa_i=1$~\cite{david2012}.

\emph{Schemes and running.} A coupling or mass in quantum field theory is defined only together with a \emph{renormalisation scheme} and
an energy scale $\mu$. In the $\overline{\text{MS}}$ scheme only the divergent part of each loop is subtracted, and parameters run with
$\mu$. The \emph{on-shell} scheme instead fixes masses at the poles of the propagators and $\alpha$ at zero momentum transfer. The same
quantity can differ between the two by more than its measurement error, so every comparison states its scheme~\cite{erler2017,navas2024}.
The dependence on $\mu$ obeys the renormalisation-group equations (RGEs). At one loop
$\dd\alpha_i^{-1}/\dd\ln\mu=-b_i/2\pi$, with $(b_1,b_2,b_3)=(41/10,-19/6,-7)$ in the Standard Model (hypercharge normalised as in grand
unification). The strong coupling therefore weakens at high energy, and the hypercharge coupling grows. Precision work runs at two or
three loops~\cite{buttazzo2013}.

\emph{Dirac or Majorana.} A \emph{Dirac} mass joins a left-handed neutrino to a right-handed partner and conserves lepton number $L$. A
\emph{Majorana} mass joins the neutrino to its own antiparticle and changes $L$ by two. The lowest-dimension operator of Standard Model
fields that gives neutrinos mass is Weinberg's dimension-five operator, $(c/\Lambda)\,(LH)(LH)$, with $L$ the lepton doublet; in the Higgs
vacuum it gives a Majorana mass of order $cv^2/\Lambda$, small when the new scale $\Lambda$ is large~\cite{weinberg1979}. Majorana masses
can therefore also come from the Higgs vacuum.

\emph{Neutrinoless double-beta decay} ($0\nu\beta\beta$) turns two neutrons of a nucleus into two protons and two electrons with no
neutrinos. It changes $L$ by two and can occur only if neutrinos have a Majorana mass. Through light-neutrino exchange its rate is
proportional to $m_{\beta\beta}^2$, with the effective mass $m_{\beta\beta}=|\sum_iU_{ei}^2m_i|$ set by the mixing matrix $U$ and the
masses $m_i$~\cite{navas2024}. Xenon and germanium experiments search for it~\cite{kamlandzen2025,legend2025}.

\emph{Grand unification.} A grand unified theory (GUT) embeds the three gauge groups in one simple group, the smallest being
$SU(5)$~\cite{georgiglashow}, with a single coupling above a scale $M_X$. With the Standard Model's particles alone, the three running
couplings do not meet at one point~\cite{amaldi}. A GUT has two generic consequences. Its heavy gauge bosons let protons decay, with a
lifetime of order $M_X^4/(\alpha_G^2m_p^5)$; Super-Kamiokande's bound on $p\to e^+\pi^0$ excludes the minimal $SU(5)$
model~\cite{superk}. And its breaking leaves magnetic monopoles of mass near $M_X/\alpha_G$, about one per causal volume, which would
overfill the universe unless inflation dilutes them~\cite{preskill,guth}.
\end{primer}

\begin{primer}{Particle physics}{box:particle}
\setlength{\parskip}{3pt plus 1pt}%
\emph{The Higgs mechanism.} The Standard Model is a gauge theory of $SU(3)\times SU(2)\times U(1)$: colour for the strong force, weak isospin
and hypercharge for the electroweak force~\cite{glashow,weinberg67,salam}. The Higgs field is a doublet whose potential has its minimum
away from zero, at a vacuum value $v=246$\,GeV~\cite{higgs64,be64}. In that vacuum the W and Z bosons acquire masses proportional to $v$,
while one combination of the electroweak fields, the photon, stays massless. Quarks and charged leptons get their masses from
\emph{Yukawa couplings} $y_f$ to the doublet, $m_f=y_fv/\sqrt2$. The Higgs boson therefore couples to each fermion with strength $m_f/v$
and to the W and Z with $2m_V^2/v$.

\emph{The $\kappa$ framework.} Measurements of Higgs production and decay are summarised by \emph{coupling modifiers} $\kappa_i$, which
scale each Standard Model coupling, so that a rate involving particle $i$ scales as $\kappa_i^2$; the Standard Model has every
$\kappa_i=1$~\cite{david2012}.

\emph{Schemes and running.} A coupling or mass in quantum field theory is defined only together with a \emph{renormalisation scheme} and
an energy scale $\mu$. In the $\overline{\text{MS}}$ scheme only the divergent part of each loop is subtracted, and parameters run with
$\mu$. The \emph{on-shell} scheme instead fixes masses at the poles of the propagators and $\alpha$ at zero momentum transfer. The same
quantity can differ between the two by more than its measurement error, so every comparison states its scheme~\cite{erler2017,navas2024}.
The dependence on $\mu$ obeys the renormalisation-group equations (RGEs). At one loop
$\dd\alpha_i^{-1}/\dd\ln\mu=-b_i/2\pi$, with $(b_1,b_2,b_3)=(41/10,-19/6,-7)$ in the Standard Model (hypercharge normalised as in grand
unification). The strong coupling therefore weakens at high energy, and the hypercharge coupling grows. Precision work runs at two or
three loops~\cite{buttazzo2013}.

\emph{Dirac or Majorana.} A \emph{Dirac} mass joins a left-handed neutrino to a right-handed partner and conserves lepton number $L$. A
\emph{Majorana} mass joins the neutrino to its own antiparticle and changes $L$ by two. The lowest-dimension operator of Standard Model
fields that gives neutrinos mass is Weinberg's dimension-five operator, $(c/\Lambda)\,(LH)(LH)$, with $L$ the lepton doublet; in the Higgs
vacuum it gives a Majorana mass of order $cv^2/\Lambda$, small when the new scale $\Lambda$ is large~\cite{weinberg1979}. Majorana masses
can therefore also come from the Higgs vacuum.

\emph{Neutrinoless double-beta decay} ($0\nu\beta\beta$) turns two neutrons of a nucleus into two protons and two electrons with no
neutrinos. It changes $L$ by two and can occur only if neutrinos have a Majorana mass. Through light-neutrino exchange its rate is
proportional to $m_{\beta\beta}^2$, with the effective mass $m_{\beta\beta}=|\sum_iU_{ei}^2m_i|$ set by the mixing matrix $U$ and the
masses $m_i$~\cite{navas2024}. Xenon and germanium experiments search for it~\cite{kamlandzen2025,legend2025}.

\emph{Grand unification.} A grand unified theory (GUT) embeds the three gauge groups in one simple group, the smallest being
$SU(5)$~\cite{georgiglashow}, with a single coupling above a scale $M_X$. With the Standard Model's particles alone, the three running
couplings do not meet at one point~\cite{amaldi}. A GUT has two generic consequences. Its heavy gauge bosons let protons decay, with a
lifetime of order $M_X^4/(\alpha_G^2m_p^5)$; Super-Kamiokande's bound on $p\to e^+\pi^0$ excludes the minimal $SU(5)$
model~\cite{superk}. And its breaking leaves magnetic monopoles of mass near $M_X/\alpha_G$, about one per causal volume, which would
overfill the universe unless inflation dilutes them~\cite{preskill,guth}.
\end{primer}
